\pdfoutput=1
% Options for packages loaded elsewhere
\PassOptionsToPackage{unicode}{hyperref}
\PassOptionsToPackage{hyphens}{url}
\documentclass[
  11pt,
]{article}
\usepackage{xcolor}
\usepackage[margin=1in]{geometry}
\usepackage{amsmath,amssymb}
\setcounter{secnumdepth}{-\maxdimen} % remove section numbering
\usepackage{iftex}
\ifPDFTeX
  \usepackage[T1]{fontenc}
  \usepackage[utf8]{inputenc}
  \usepackage{textcomp} % provide euro and other symbols
\else % if luatex or xetex
  \usepackage{unicode-math} % this also loads fontspec
  \defaultfontfeatures{Scale=MatchLowercase}
  \defaultfontfeatures[\rmfamily]{Ligatures=TeX,Scale=1}
\fi
\usepackage{lmodern}
\ifPDFTeX\else
  % xetex/luatex font selection
\fi
% Use upquote if available, for straight quotes in verbatim environments
\IfFileExists{upquote.sty}{\usepackage{upquote}}{}
\IfFileExists{microtype.sty}{% use microtype if available
  \usepackage[]{microtype}
  \UseMicrotypeSet[protrusion]{basicmath} % disable protrusion for tt fonts
}{}
\makeatletter
\@ifundefined{KOMAClassName}{% if non-KOMA class
  \IfFileExists{parskip.sty}{%
    \usepackage{parskip}
  }{% else
    \setlength{\parindent}{0pt}
    \setlength{\parskip}{6pt plus 2pt minus 1pt}}
}{% if KOMA class
  \KOMAoptions{parskip=half}}
\makeatother
\setlength{\emergencystretch}{3em} % prevent overfull lines
\providecommand{\tightlist}{%
  \setlength{\itemsep}{0pt}\setlength{\parskip}{0pt}}
\usepackage[]{natbib}
\bibliographystyle{plainnat}
\usepackage{bookmark}
\IfFileExists{xurl.sty}{\usepackage{xurl}}{} % add URL line breaks if available
\urlstyle{same}
\hypersetup{
  pdftitle={Example 3 - Dentist},
  hidelinks,
  pdfcreator={LaTeX via pandoc}}

\title{Example 3 - Dentist}
\author{}
\date{May 2025}

\begin{document}
\maketitle

\subsection{Dentist}\label{dentist}

\textbf{1. Job Composition}

\begin{itemize}
\tightlist
\item
  Thinking: 25\% (diagnosis, treatment planning, medical knowledge)
\item
  Physical: 30\% (precise hand movements, tool manipulation)
\item
  Social: 15\% (patient communication, anxiety management)
\item
  Perception: 20\% (visual examination, tactile feedback)
\item
  Adaptation: 10\% (responding to patient movement, unexpected findings)
\end{itemize}

\textbf{2. Human Advantage Score}

\begin{itemize}
\tightlist
\item
  Thinking: 5/10 (AI can assist with diagnosis but struggles with
  comprehensive treatment planning)
\item
  Physical: 9/10 (extremely difficult to replicate precise hand
  movements in small, sensitive spaces)
\item
  Social: 7/10 (difficult to replicate bedside manner and real-time
  emotional reassurance)
\item
  Perception: 8/10 (complex integration of visual and tactile feedback
  in variable environments)
\item
  Adaptation: 9/10 (significant challenges adapting to unexpected
  patient movements and oral conditions)
\end{itemize}

\textbf{3. Job Resilience Score}

\begin{itemize}
\tightlist
\item
  Total: 7.5/10
\item
  Beginner: 6.5 (7.5-1)
\item
  Advanced: 7.5
\item
  Master: 8.5 (7.5+1)
\end{itemize}

\textbf{4. Result:}

\begin{itemize}
\tightlist
\item
  Beginner: Low Risk (10+ years)
\item
  Advanced: Low Risk (10+ years)
\item
  Master: Very Low Risk (15+ years or indefinite human advantage)
\end{itemize}

\textbf{5. Adoption Factors Check:} High salaries might justify
automation investment (risk up), but implementation costs for robotic
dentistry would be extremely high, and many patients are likely to
prefer humans (risk down). These factors cancel each other out, keeping
dentistry in the Low to Very Low Risk category.

\end{document}
